\section{Discussion}
\label{sec:discussion}

\paragraph{Implications for audit design.} Audit design faces a failure mode distinct from threshold miscalibration: a well-calibrated threshold can still produce worse long-run outcomes than no safe harbor at all, if certification confers a sufficiently strong one, structurally rather than as a one-time artifact (Theorem~\ref{prop:recurring}). Case (a)'s degraded steady state is invisible from pass/fail records, since such a firm never subsequently fails; Case (b)'s cycling surfaces as repeated failing and re-passing, but nothing flags why. Both underperform the no-safe-harbor counterfactual, and Section~\ref{sec:sim-heterogeneous} puts a number on it: roughly a third of achievable quality under a stylized population. One lever independent of the threshold: raising $\rho$ toward $1$ (e.g., preserving contestability rights after certification). A second, more speculative lever — shortening the audit interval relative to drift — is suggestive but not evaluable here, since our model normalizes cadence to a single period; Ball and Knoepfle~\cite{ball2023timing} give the closer formal treatment.

\paragraph{Imperfect observation is a third lever, with a caveat.} Proposition~\ref{prop:imperfectobs} adds a lever the first two do not suggest: a noisily-observed threshold gives a single crossing point in place of two discrete levels, and above threshold $\sigma_0$ this point is stable and, in our simulations, raises quality — against the usual presumption that noisier audits are worse. But below $\sigma_0$ the point is unstable and cycling persists instead, converging to the hard-threshold model's behavior as noise vanishes. What matters is not imprecision per se but removing the cliff edge a firm can fully escape past; whether that stabilizes or destabilizes depends on how fast scrutiny changes near the crossing point, which is why the benefit reverses at small $\sigma$. A noisy audit that is also biased or gameable is not covered. We read this as a reason to distinguish cliff-edge-removing noise from noise that merely degrades information, not a recommendation to make audits noisier on purpose.

\paragraph{The regulator's design problem, and what remains open.} Welfare is non-decreasing in $\tau$ up to $q^*(p_H)$ and flat thereafter, but $\tau=q^*(p_H)$ attains that maximum only by making certification unattainable, so the recommendation is $\tau$ close to, but strictly below it — a numerically-supported conjecture, holding $\rho$ fixed (Section~\ref{sec:regulator}). This falls short of the fully strategic problem: jointly choosing $\tau$ and a policy shaping $\rho$ under heterogeneity (Section~\ref{sec:heterogeneous}) rather than a representative firm. Section~\ref{sec:related}'s mechanism-design literature~\cite{machostadler2006optimal,huang2024zetar,varas2020random} gives relevant machinery for a related but distinct problem; combining it with our results under heterogeneity is the natural next extension.

\paragraph{Limitations.} This is primarily theoretical: Sections~\ref{sec:sim}--\ref{sec:sim-functionalform} show the central results are satisfiable and not an artifact of functional form, but the imperfect-observation and heterogeneity results are checked only at our calibration. Section~\ref{sec:calibration} anchors $\tau=0.80$ against real disclosure statistics but cannot test the dynamic predictions directly, since the needed firm-level panel does not yet exist. The logistic detection function is a modeling convenience, and the noise-raises-quality finding holds only above $\sigma_0$, established numerically. Section~\ref{sec:regulator}'s result likewise holds $\rho$ fixed, with a single representative firm.
